% commentaire.tex — f164v, vv. 165–207 (conversion initiale des travaux antérieurs)
\comm{168–70}{Φιλοκλῆς, Ξενοκλῆς, Θέογνις.}{Trois poètereaux tragiques contemporains d'Aristophane.\\
À propos du vilain Philoclès, deux scholies nous ont été transmises dans le manuscrit de Ravenne\nc{Manuscrit qui possède l'unique exemplaire des \italique{Thesmophories} qui nous soit parvenu -- \italique{unicodex} dit-on. Le manuscrit de Munich (\nieme{XV} siècle) n'étant qu'une recopie du premier (voir l'édition de von Velsen des \italique{Aristophanis Thesmophoriazusae} (1883) pour une comparaison complète du \italique{Ravennas} (R) et du \italique{Monacensis} (G)).} :\citComm{Φιλοκλέης: τραγῳδίας ποιητὴς, αἰσχρὸς τὸ πρόσωπον καὶ σιμὸς καὶ μικρὸς τὸ σῶμα\nc{Suidas (φ 379) reprend cette scholie qu'il joint à d'autres éléments tirés du scholiaste aux \italique{Oiseaux} (v. 281).}.\\
ἐκ τούτου δηλονότι καὶ ἐν τοῖς Ὄρνισιν (v. 1295) οὐκ ἄλλως κορυδῷ αὐτὸν ὡ[μοίου] ἀλλ᾽ ἢ ὅτι μικρὸς καὶ αἰσχρός. καὶ Τηλεκλείδης Ἡσιόδοις\nc{fr. 15 in \italique{Poetae Comici Graeci}, VII, edd. Kassel et Austin.}·
`ἀλλ᾽ ἡ τάλαινα Φιλοκλέα βδε\douteux{λ}\douteux{λ}...\douteux{λ}.οθεν οὖν·
εἰ δ᾽ ἐστὶν Αἰσχύλου φρόνημ᾽ ἔχων'.}Le scholiaste ne se trompe pas lorsqu'il prétend que la comparaison avec l'\italique{alouette huppée}, κορυδός, était d'ordre physique\nc{Ce qui explique le surnom du \italique{Coryde} dont il est question dans les \italique{Deipnosophistes} d'Athénée (VI, 241a \&c.) «~qui aimait à se faire persifler~».} : on connaît par les parémiographes\nc{Zénobius IV, 59 ; Diogenianus V, 57 ; Michel Apostolius IX, 92 ; \&c. à retrouver dans le \italique{Corpus paroemiographorum graecorum}, ed. Leutsch. Cf. aussi Hésychius (σ 3681). Kassel \& Austin l'ont inclus (n° 827) aux \italique{Comica Adespota} de leur édition des \italique{Poetae Comici Graeci}, VIII.} le proverbe grec «~être plus hideux que Corydée~» :\citComm{Κορυδέως εἰδεχθέστερος· ἐπὶ δυσμορφίᾳ ἐκωμῳδεῖτο καὶ οὗτος καὶ οἱ παῖδες αὐτοῦ ὧν ἕνα τὸν Ἀρχέστρατον φασιν.}La pauvre alouette huppée n'est pas bien lotie par la tradition populaire : en plus d'être laide, même dans la force de l'âge elle est chétive (plus que l'aigle vieillissant : ἀετοῦ γῆρας, κορύδου νεότης), et n'a pas pour autant le ramage du rossignol (ἐν ἀμούσοις καὶ κόρυδος φθέγγεται\nc{Voir ce qu'en dit Érasme, qui a recensé ce proverbe dans ses \italique{Adages} (n° 1192).})... « au pays des aveugles, les borgnes sont rois ».\\
(Voir le commentaire au vers 441 pour le cas Xénoclès.)\\
Jean Taillardat a consacré la note que voici à Théognis\nc{\italique{Les images d'Aristophane} (1965), §759.} :\\
Théognis le Tragique est ennuyeux et \italique{froid}, ψυχρός, au point qu'on le surnomme la \italique{Neige} (Χιών, Σ \italique{Acharniens} 11). La représentation d'une de ses pièces suffit à déchaîner les rigueurs de l'hiver, même sous le ciel déjà peu clément de la Thrace (\italique{Ach.} 138-140). Pris de pitié pour un froid rimailleur, Pisthétaire lui offre un vêtement de peau (\italique{Oiseaux} 935) : 
\citComm{ἔχε τὴν σπολάδα· πάντως δέ μοι ῥιγῶν δοκεῖς.\\
\francais{Prends ce manteau ; tu me parais complètement gelé.}}
Cette plaisanterie est continuée aux vers 946 et 955 (ταυταγὶ | τὰ κρυερὰ, \italique{ces vers glacés} ; cf. νιφοβόλους ἀναβολάς (\italique{ib.} 1385), « des préludes battus de neige ». Cet emploi de ψυχρός est banal\nc{Ψυχρός en parlant d'un \italique{poète} : Alexis 179 [$=$ 184 Kassel-Austin $=$ Athénée, III, 123e], Machon (\italique{ap.} Athénée 580a), d'un \italique{orateur} : Plat. \italique{Euthyd.} 284e ; Théophil. 4, 4, \&c. Ψυχρός se dit aussi d'une tragédie manquée (\italique{Thesm.} 848), d'une plaisanterie sans sel (Eupol. 244 [$=$ 261 Kassel-Austin]).}.)}
\comm{172–3}{πῶς πρὸς τῶν θεῶν; — παῦσαι βαΰζων.}{Goûtons à l'allitération en `t', `p' et surtout en `o' aboyée par Mnésiloque : π\textbf{ῶ}ς πρ\textbf{ὸ}ς τ\textbf{ῶ}ν θε\textbf{ῶ}ν ; c'est cette dissonance que lui reproche Euripide – pour la forme, car la question de Mnésiloque est aussi tout-à-fait désobligeante et ironique... il est certainement à l'affût de quelque double sens au «~ἐμαυτὸν ἐθεράπευσα~» d'Agathon. 
Euripide interrompt donc : παῦσαι βαΰζων, « cesse d'aboyer ! » -- quoique le verbe «~aboyer~» ne rende pas le registre tragique du verbe βαΰζω, formé par onomatopée sur le cri βαῢ βαΰ du chien (αὗ αὗ au vers 903 des \italique{Guêpes}).}
\comm{173–5}{καὶ γὰρ ἐγὼ τοιοῦτος ἦν ... — οὐ ζηλῶ σε τῆς παιδεύσεως.}{Euripide se reconnaît – non sans quelque condescendance – une parenté artistique\nc{Confirmée par Aristote. (Voir le commentaire \italique{infra} au vers 29, n. 3.)} quand Mnésiloque répond comme si Agathon était héritier des mœurs d'Euripide.}
\comm{177–8}{Ἀγάθων σοφοῦ πρὸς ἀνδρός ὅστις ἐν βραχεῖ / πολλοὺς καλῶς οἷός τε συντέμνειν λόγους.}{Citation quasi-textuelle de l'\italique{Éole} (Αἰόλος) d'Euripide rapportée par Stobée (3, 35, 3)\nc{À retrouver dans l'édition de C. Wachsmuth \& O. Hense, \italique{Ioannis Stobaei Anthologium}, vol. 3, 1904, p. 688 l. 5-7. C'est le fragment 21 de l'édition Jouan-van Looy aux Belles Lettres, t. VIII, 1\italique{ère} partie ; il correspond respectivement aux fragments 28 et 33 des éditions Kannicht et Mette.} :\citComm{παῖδες, σοφοῦ πρὸς ἀνδρός ὅστις ἐν βραχεῖ,\\
πολλοὺς λόγους οἷός τε συντέμνειν καλῶς}P.-L. Courier, \italique{Pamphlet des pamphlets} :
\citComm{\francais{C'est un mérite non commun, ni facile, de clore en peu de mots beaucoup de sens.}(Note de Willems.)}}
\comm{180}{τοῦ χρείαν ἔχων;}{Euripide, \italique{Suppliantes}, 115 : Τί χρῆμα θηρῶν καὶ τίνος χρείαν ἔχων;}
\comm{186}{σαφῶς σώσεις ἐμέ.}{Le sigmatisme tant reproché à Euripide ! toujours avec le verbe σῴζειν, \italique{sauver}, le véritable Euripide avait sifflé au vers 476 de sa \italique{Médée} : ἔσῳσά σ' ὡς ἴσασιν Ἑλλήνων ὅσοι. «~Grand merci de nous avoir sauvés des sigmas d'Euripide~» avait persiflé Platon le Comique\nc{Fragment 29(30) de l'édition Kassel-Austin, \italique{Poetae Comici Graeci}, vol. VII (1989) ; chez Meineke, \italique{Fragmenta Comicorum Graecorum}, il se trouve au tome II, p. 626.} en reprenant le début de vers : \citComm{εὖ γέ σοι γένοιθ᾽ ἡμᾶς ὅτι\\
ἔσωσας ἐκ τῶν σίγμα τῶν Εὐριπίδου.}}
\comm{193}{χαίρεις ὁρῶν φῶς, πατέρα δ᾽ οὐ χαίρειν δοκεῖς;}{Citation tirée de l'\italique{Alceste} d'Euripide, v. 691 : Phérès envoie paître son fils Admète qui lui demande de mourir à sa place. C'est finalement Alceste qui se sacrifiera pour son mari ; tout comme Mnésiloque, parent par \italique{alliance} d'Euripide, ne va pas tarder à accepter ce qui ressemble fort à un suicide.}
\comm{200–1}{εὐρύπρωκτος εἶ / οὐ τοῖς λόγοισιν ἀλλὰ τοῖς παθήμασιν.}{Mnésiloque reprend l'opposition τοῖς τεχνάσμασιν / τοῖς παθήμασιν d'Agathon en substituant τοῖς λόγοισιν à τοῖς τεχνάσμασιν.\\
Le spectateur attend donc l'opposition classique τοῖς λόγοις / τοῖς ἔργοις : \italique{en théorie} / \italique{en pratique}... sauf qu'ici les \italique{pratiques} d'Agathon sont jugées par trop passives.}
\comm{204}{νυκτερείσια.}{Pour le texte, voir le commentaire critique.}
\comm{205}{ὑφαρπάζειν τε θήλειαν Κύπριν.}{L'expression ὑφαρπάζειν Κύπριν se retrouve au \nieme{722} vers de l'\italique{Assemblée des femmes}.
Les Athéniennes ont toujours une bonne raison d'en vouloir aux poètes tragiques : lorsqu'ils ne médisent pas d'elles, ils se livrent à une concurrence déloyale.\citComm{Mais les jeunes femmes d'Athènes\\
Privées d'leurs époux, d'leurs galants,\\
Accumulèrent la rancune patiemment...\\
Puis un jour, ivres de colère,\\
Elles s'armèrent de bâtons,\\
Et, farouches, elles immolèrent Agathon !}}
\comm{207}{ἀτὰρ ἡ πρόφασίς γε νὴ Δί᾽ εἰκότως ἔχει.}{Molière, \italique{Misanthrope}, IV, 3 :
\citComm{\francais{Ah! le détour est bon et l'excuse admirable.}}\textbf{La défaite}: prétexte fallacieux pour se défaire de ce qui est importun, embarrassant. (J'ai laissé le mot soigneusement choisi par Willems.)}
